\chapter{The pre-transformer era}
\label{ch:pre-transformer}

\newthought{The transformer is a 2017 paper.} Language modelling has been a
serious research topic for at least 70 years before that. Most of the
intuitions in modern LLMs — the chain-rule decomposition, the importance of
context, the trade-off between fixed and unbounded memory — are not new. They
have been re-discovered, with increasing computational firepower, several
times.

This chapter is a tour of what came before the transformer, framed in the
language we will use in the rest of the book. It is short, but it matters:
many of the ``new'' subquadratic architectures in Part~V are direct
descendants of pre-transformer ideas that the field stopped using because
GPUs got big enough to brute-force the problem.

\section{n-gram models — counting your way to fluency}

The first practical language models were just \emph{count tables}. Pick a
window size \(k\) (usually 3 to 5). For every \(k\)-tuple of tokens you ever
saw in a training corpus, store
\begin{equation}
  p(x_t \mid x_{t-k+1}, \ldots, x_{t-1})
  \;=\;
  \frac{\#(x_{t-k+1}, \ldots, x_{t-1}, x_t)}
       {\#(x_{t-k+1}, \ldots, x_{t-1})}.
  \label{eq:ngram}
\end{equation}
That is the entire model.\sidenote{In practice you smooth the counts so
unseen tuples get a tiny non-zero probability — Kneser-Ney smoothing
\citep{kneser1995smoothing} was the dominant technique for two decades.}

\paragraph{Strengths.} O(1) compute per token. Embarrassingly parallel to
train. For \(k = 5\), competitive with much fancier models on speech-recognition
language-model rescoring up to roughly 2010.\sidenote{Google's 2007
paper described training 5-grams over 2 trillion words for machine
translation, several years before any neural alternative was viable
\citep{brants2007large}.}

\paragraph{Weaknesses.} Fundamentally bounded by \(k\). An \(n\)-gram model
literally cannot use information from more than \(k - 1\) tokens back, even
in principle. The state space grows as \(V^k\), so going from 5-grams to
10-grams blows up by twelve orders of magnitude.

That hard bound — \emph{context up to \(k\) tokens, period} — is the original
sin every architecture since has been trying to escape.

\section{Recurrent neural networks (RNNs)}

A recurrent network throws the count table out and replaces it with a single
hidden vector \(h_t \in \R^{d}\) that is updated step by step:
\begin{equation}
  h_t \;=\; \tanh(W_h h_{t-1} + W_x e_t + b),
  \quad
  p_\theta(x_{t+1} \mid x_{\le t}) \;=\; \softmax(W_o h_t).
  \label{eq:rnn}
\end{equation}
Now the model can in principle use \emph{all} prior tokens, because each
\(h_t\) is a function of the entire prefix.\sidenote{The clean derivation
goes back to \citep{elman1990finding}; \citep{mikolov2010rnnlm} put it on
the LM map for modern NLP.} Compute per token is \(\bigO(d^2)\), independent
of \(t\), which is a beautiful property. We will return to that.

\paragraph{The vanishing gradient problem.} The catch is that ``in
principle'' is doing all the work. Training the RNN by backpropagation
through time multiplies many Jacobians together; with the \(\tanh\)
non-linearity, the gradients vanish (or, worse, explode) within roughly
20 steps. Vanilla RNNs effectively have a context of about that length, and
nobody could train them stably enough to do better.

\section{LSTMs and GRUs — the gated decade}

The fix that worked is \emph{gating}. A long short-term memory cell
\citep{hochreiter1997lstm} replaces the simple update in
equation~\eqref{eq:rnn} with a small machine made of \term{gates}:
\begin{align}
  f_t &= \sigma(W_f [h_{t-1}, e_t] + b_f) & \text{forget gate} \\
  i_t &= \sigma(W_i [h_{t-1}, e_t] + b_i) & \text{input gate} \\
  o_t &= \sigma(W_o [h_{t-1}, e_t] + b_o) & \text{output gate} \\
  \tilde c_t &= \tanh(W_c [h_{t-1}, e_t] + b_c) & \text{candidate} \\
  c_t &= f_t \odot c_{t-1} + i_t \odot \tilde c_t & \text{cell state} \\
  h_t &= o_t \odot \tanh(c_t).
\end{align}
The cell state \(c_t\) flows along an additive path, so its gradient does
not vanish; the gates decide what to write, what to keep, and what to read.
\citet{cho2014gru} simplified this to the GRU, which has fewer parameters and
behaves similarly in practice.\sidenote{If you have ever wondered why
\emph{Mamba} (Chapter~\ref{ch:ssm}) keeps using the words ``selective'',
``state'', and ``gate'' — it is because Mamba is the same idea, polished
and re-discovered with hardware-aware tricks.}

LSTMs ran the world from roughly 2014 to 2017. Google Translate, Siri's
on-device speech, the original Alexa wake-word detector — all LSTM under
the hood.

\paragraph{Why LSTMs eventually lost.} Two reasons.

\begin{enumerate}
\item \textbf{Sequential by construction.} Each \(h_t\) depends on \(h_{t-1}\).
You cannot parallelise across the time axis; the GPU sits half idle.
Transformers, in contrast, evaluate every position simultaneously during
training.
\item \textbf{Compressed memory.} Everything from the past has to be squeezed
into a single \(h_t\) vector of fixed size. With enough capacity an LSTM
can remember a few thousand tokens; beyond that, things blur. We will see
this exact trade-off come back as a feature of state-space models, where it
is sometimes a virtue.
\end{enumerate}

\section{Encoder-decoder, sequence-to-sequence, and the original ``attention''}

In 2014, two papers \citep{sutskever2014seq2seq, cho2014gru} proposed
\term{sequence-to-sequence} models for machine translation: an encoder LSTM
reads the source sentence into a fixed-size vector, a decoder LSTM generates
the target one token at a time, conditioned on that vector.

This worked for short sentences and fell over for long ones — the bottleneck
of compressing an entire paragraph into one vector was too tight.

The fix is the original meaning of the word \term{attention}.
\citet{bahdanau2014attention} let the decoder, at every step, look back at
\emph{all} the encoder hidden states and produce a weighted average:
\begin{equation}
  c_t \;=\; \sum_{i=1}^{n} \alpha_{t,i}\, h_i^{\text{enc}},
  \quad
  \alpha_{t,i} \;=\; \frac{\exp(\text{score}(h_t^{\text{dec}}, h_i^{\text{enc}}))}
                          {\sum_{j} \exp(\text{score}(h_t^{\text{dec}}, h_j^{\text{enc}}))}.
\end{equation}
The weights \(\alpha_{t,i}\) say ``how much should the decoder pay attention
to source position \(i\) when emitting target token \(t\)''. They are softmax
normalised. They are computed from the current decoder state and the encoder
states.

\keyidea{The attention layer that runs every transformer in 2026 is a direct
descendant of the Bahdanau attention bolted onto an LSTM in 2014. The
transformer's contribution is to delete the LSTM and use \emph{only}
attention.}

\section{The 2017 turn}

What \citet{vaswani2017attention} pointed out is that you can build a
sequence model out of attention alone. No recurrence, no convolution.
Stack attention layers, mix in some position-wise feed-forward layers, add
positional encodings so the model knows token order, and you get a model
that:

\begin{itemize}
\item parallelises across the time axis (every position is computed at
once),
\item lets every token attend to every other token in a single layer (no
vanishing gradient over distance),
\item and trains roughly an order of magnitude faster on the same hardware.
\end{itemize}

The price is a quadratic dependence on sequence length. That price is what
the rest of the book is about.

The 2017 paper was about machine translation, but within four years the same
recipe — scaled up by three orders of magnitude — was producing GPT-3, a
language model that could do arithmetic, write code, and pass standardised
exams. The transformer ate the field. Part~II is the careful explanation
of why.

\fillme{FILL-03-01}{Code listing: vanilla RNN forward pass in PyTorch}{%
  ${\sim}20$-line PyTorch listing implementing equation~\ref{eq:rnn} as a
  loop over time, returning the per-step hidden states. Use the
  \texttt{lstlisting} environment. Include shape comments on every line that
  manipulates a tensor. Below the listing, add a one-paragraph note that
  this is correct but slow because of the time-axis loop, foreshadowing the
  parallelism advantage of attention in Chapter~\ref{ch:transformer}.}{%
  PyTorch \texttt{nn.RNNCell} source as a reference; equation~\ref{eq:rnn}.}{%
  ${\sim}20$ lines of code + one paragraph.}

\fillme{FILL-03-02}{Code listing: LSTM cell forward pass in PyTorch}{%
  Same shape as FILL-03-01, but for a single LSTM cell. Render the four
  gates explicitly so the reader can map every line to one row of the
  derivation in this chapter. Comment which lines correspond to ``forget'',
  ``input'', ``output'', and ``cell update''. Don't use the built-in
  \texttt{nn.LSTMCell} — write the linear algebra by hand.}{%
  Hochreiter \& Schmidhuber 1997; PyTorch source for sanity-check.}{%
  ${\sim}25$ lines of code + 2 sentence caption.}

\fillme{FILL-03-03}{Worked example: parameter count of an LSTM cell}{%
  Derive symbolically the parameter count of an LSTM cell with input
  dimension \(d\) and hidden dimension \(h\). Show the answer is
  \(4 \cdot (d \cdot h + h \cdot h + h)\). Plug in \(d = h = 1024\) for a
  numeric. Then compare to a transformer block of comparable hidden
  dimension (forward-pointer to Chapter~\ref{ch:transformer}).}{%
  Standard LSTM derivation; equation block from this chapter.}{%
  Half a page; one symbolic equation, one numeric.}
